\documentclass[11pt,a4paper]{article}

% ── Codificación, fuentes y lenguaje ───────────────────────────────────────────
\usepackage[utf8]{inputenc}
\usepackage[T1]{fontenc}
\usepackage[default,scale=0.95]{sourcesanspro}
\renewcommand{\familydefault}{\sfdefault}
\usepackage[spanish,es-noshorthands,es-tabla]{babel}

% ── Geometría y espaciado ─────────────────────────────────────────────────────
\usepackage[top=2.2cm, bottom=2.5cm, left=2.2cm, right=2.2cm, headheight=15pt]{geometry}
\usepackage{setspace}
\setstretch{1.18}
\usepackage{parskip}

% ── Matemáticas ───────────────────────────────────────────────────────────────
\usepackage{amsmath}
\usepackage{amssymb}

% ── TikZ (simbolos ISO 5807 / ANSI X3.5) ──────────────────────────────────────
\usepackage{tikz}
\usetikzlibrary{babel}
\usetikzlibrary{arrows.meta, positioning, shapes.geometric, shapes.symbols,
                decorations.pathmorphing, calc}

% ── Tablas y listas ───────────────────────────────────────────────────────────
\usepackage{booktabs}
\usepackage{tabularx}
\usepackage{array}
\usepackage{enumitem}

% ── Colores, cajas e iconos ───────────────────────────────────────────────────
\usepackage[dvipsnames]{xcolor}
\usepackage{tcolorbox}
\tcbuselibrary{skins, breakable}
\usepackage{fontawesome5}

% ── Secciones con barra de color (infografia) ─────────────────────────────────
\usepackage{titlesec}
\titleformat{\section}
  {\normalfont\LARGE\bfseries\color{azulNoche}}
  {\colorbox{cyanNeon}{\color{fondoOscuro}\thesection}}{0.7em}{}
  [\vspace{2pt}\color{cyanNeon}\rule{\linewidth}{1.8pt}]
\titlespacing*{\section}{0pt}{24pt}{10pt}
\titleformat{\subsection}
  {\normalfont\large\bfseries\color{azulNoche}}
  {\textcolor{cyanNeon}{\thesubsection}}{0.7em}{}
  [\vspace{1pt}\color{grisLinea}\rule{\linewidth}{0.6pt}]
\titlespacing*{\subsection}{0pt}{14pt}{6pt}

% ── Cabeceras y pies de página ────────────────────────────────────────────────
\usepackage{fancyhdr}
\usepackage{lastpage}
\pagestyle{fancy}
\fancyhf{}
\renewcommand{\headrulewidth}{0pt}
\renewcommand{\footrulewidth}{0pt}
\fancyfoot[C]{%
  \begin{minipage}{\linewidth}
    \vspace{2pt}
    \color{cyanNeon}\rule{\linewidth}{1.2pt}\\[2pt]
    \centering\color{grisTexto}\footnotesize
    Página \thepage\ de \pageref{LastPage}
  \end{minipage}%
}

% ── Hiperenlaces ──────────────────────────────────────────────────────────────
\usepackage[colorlinks=true, linkcolor=azulNoche, urlcolor=cyanNeon]{hyperref}
\sloppy
\emergencystretch 3em

% ── Paleta de colores ─────────────────────────────────────────────────────────
\definecolor{fondoOscuro}{HTML}{0D1B2A}
\definecolor{verdeTurquesa}{HTML}{004D40}
\definecolor{azulNoche}{HTML}{1B3A6B}
\definecolor{azulMedio}{HTML}{2A5298}
\definecolor{cyanNeon}{HTML}{00C8E0}
\definecolor{verdeSignal}{HTML}{00B887}
\definecolor{naranjaVivo}{HTML}{FF7043}
\definecolor{rojoAlerta}{HTML}{E53935}
\definecolor{grisPapel}{HTML}{F4F6FA}
\definecolor{grisLinea}{HTML}{C8D0E0}
\definecolor{grisTexto}{HTML}{6B7A99}
\definecolor{amarilloNota}{HTML}{FFB300}

% ── Banda de título ───────────────────────────────────────────────────────────
\newcommand{\iconoBanda}{project-diagram}
\newcommand{\bandaTitulo}[3][fondoOscuro]{%
  \begin{tcolorbox}[
    enhanced,
    colback=#1, colframe=#1,
    boxrule=0pt, arc=8pt, outer arc=8pt,
    width=\linewidth,
    top=18pt, bottom=6pt, left=14pt, right=14pt,
    borderline south={3.5pt}{0pt}{cyanNeon},
  ]
    \begin{minipage}[c]{0.07\linewidth}
      \centering
      \textcolor{cyanNeon}{\fontsize{30}{30}\selectfont\faIcon{\iconoBanda}}
    \end{minipage}%
    \hspace{8pt}%
    \begin{minipage}[c]{0.88\linewidth}
      {\fontsize{20}{24}\selectfont\bfseries\color{white}#2}\par\vspace{5pt}%
      \textcolor{cyanNeon!80!white}{\small\faIcon{angle-right}~#3}%
    \end{minipage}
    \vspace{4pt}
  \end{tcolorbox}%
  \vspace{8pt}%
}

% ── Tarjetas ──────────────────────────────────────────────────────────────────
\tcbset{
  cajaTitulo/.style={
    enhanced, breakable,
    arc=6pt, outer arc=6pt,
    colback=azulNoche!10!grisPapel, colframe=azulNoche,
    boxrule=1pt,
    fonttitle=\bfseries, coltitle=white,
    attach boxed title to top left={yshift=-3mm, xshift=6mm},
    boxed title style={colback=azulNoche, colframe=azulNoche, arc=4pt, boxrule=0pt},
    drop fuzzy shadow=azulNoche!25!white,
    top=8pt, bottom=8pt, left=10pt, right=10pt,
  },
  cajaContenido/.style={
    enhanced, breakable,
    arc=5pt, outer arc=5pt,
    colback=grisPapel, colframe=grisLinea,
    boxrule=0.8pt,
    fonttitle=\bfseries\small, coltitle=azulNoche,
    top=6pt, bottom=6pt, left=10pt, right=10pt,
  },
  cajaConcepto/.style={
    enhanced, breakable,
    arc=6pt, outer arc=6pt,
    colback=verdeSignal!8!white,
    colframe=verdeSignal, boxrule=1pt,
    borderline west={4pt}{0pt}{verdeSignal},
    fonttitle=\bfseries\small\color{white},
    title={\faIcon{lightbulb}~Concepto clave},
    attach boxed title to top left={yshift=-3mm, xshift=6mm},
    boxed title style={colback=verdeSignal!80!black, arc=4pt, boxrule=0pt},
    drop fuzzy shadow=verdeSignal!20!white,
    left=10pt, right=10pt, top=8pt, bottom=8pt,
  },
  cajaMejora/.style={
    enhanced, breakable,
    arc=6pt, outer arc=6pt,
    colback=naranjaVivo!8!white, colframe=naranjaVivo,
    boxrule=1pt,
    borderline west={4pt}{0pt}{naranjaVivo},
    fonttitle=\bfseries\small, coltitle=naranjaVivo!80!black,
    drop fuzzy shadow=naranjaVivo!20!white,
    top=6pt, bottom=6pt, left=10pt, right=10pt,
  },
}

\newtcolorbox{cajaConcepto}{
    enhanced, breakable,
    arc=6pt, outer arc=6pt,
    colback=verdeSignal!8!white,
    colframe=verdeSignal, boxrule=1pt,
    borderline west={4pt}{0pt}{verdeSignal},
    fonttitle=\bfseries\small\color{white},
    title={\faIcon{lightbulb}~Concepto clave},
    attach boxed title to top left={yshift=-3mm, xshift=6mm},
    boxed title style={colback=verdeSignal!80!black, arc=4pt, boxrule=0pt},
    drop fuzzy shadow=verdeSignal!20!white,
    left=10pt, right=10pt, top=8pt, bottom=8pt,
}

\newcommand{\iconotexto}[2]{\textcolor{cyanNeon}{\faIcon{#1}}~#2}

% ── Estilos de flujo (ISO 5807 / ANSI X3.5) ───────────────────────────────────
\tikzset{
  flujo/.style={
    >=Stealth,
    font=\footnotesize\sffamily,
  },
  terminador/.style={
    draw=azulNoche, fill=azulNoche!8!white, rounded corners=3pt,
    text width=1.6cm, align=center, inner sep=2mm,
    font=\footnotesize\sffamily,
  },
  proceso/.style={
    draw=azulNoche, fill=cyanNeon!10!white,
    text width=1.6cm, align=center, inner sep=2mm,
    font=\footnotesize\sffamily,
  },
  entradasalida/.style={
    draw=azulNoche, fill=verdeSignal!10!white,
    trapezium, trapezium left angle=75, trapezium right angle=105,
    text width=1.8cm, align=center, inner sep=2mm,
    font=\footnotesize\sffamily,
  },
  decision/.style={
    draw=azulNoche, fill=amarilloNota!20!white,
    diamond, aspect=1.6, text width=1.5cm, align=center, inner sep=1pt,
    font=\footnotesize\sffamily,
  },
  almacenamiento/.style={
    draw=azulNoche, fill=grisPapel,
    cylinder, shape border rotate=90, aspect=0.3,
    text width=1.7cm, align=center, inner sep=2mm,
    font=\footnotesize\sffamily,
  },
  documento/.style={
    draw=azulNoche, fill=azulMedio!8!white,
    text width=1.8cm, align=center, inner sep=2mm,
    font=\footnotesize\sffamily,
  },
  nota/.style={
    draw=grisLinea, dashed, fill=grisPapel,
    text width=1.8cm, align=center, inner sep=2mm,
    font=\scriptsize\sffamily,
  },
  etiqueta/.style={font=\scriptsize\sffamily, fill=white, fill opacity=0.75, text opacity=1, inner sep=1pt},
}

% Onda inferior para el simbolo de documento (ISO 5807)
\newcommand{\ondasDoc}[1]{%
  \draw[azulNoche, line width=0.6pt, decorate,
        decoration={snake, amplitude=0.8pt, segment length=2.4mm}]
    ([yshift=-2pt]#1.south west) -- ([yshift=-2pt]#1.south east);
}

% ── Cabecera específica del documento ─────────────────────────────────────────
\fancyhead[L]{\small\color{azulNoche}\textbf{ELECTRÓNICA} --- FAQ Higiene de Estado (Diagramas de Flujo)}
\fancyhead[R]{\small\color{grisTexto}11/09/2026}
\renewcommand{\headrulewidth}{0pt}

% ─────────────────────────────────────────────────────────────────────────────
\begin{document}

\bandaTitulo{FAQ Higiene de Estado de Sesión --- Diagramas de Flujo}{Traducción de \texttt{faq\_higiene\_estado\_sesion.md} a simbología ISO 5807 / ANSI X3.5: cadena de integridad, guardia MCP, veredictos de \texttt{check}/\texttt{clean} y eslabón semántico de bitácoras}

\vspace{4pt}
\begin{center}
\begin{tcolorbox}[cajaContenido, title={\faIcon{map-signs}~Leyenda de símbolos---ISO 5807 / ANSI X3.5}]
\small
Óvalo = \textbf{terminador} (inicio/fin) $\cdot$
Rectángulo = \textbf{proceso} (acción determinista) $\cdot$
Rombo = \textbf{decisión} (ramas Sí/No) $\cdot$
Paralelogramo = \textbf{entrada/salida} (lectura/escritura de archivos) $\cdot$
Cilindro = \textbf{almacenamiento online} (\texttt{session\_log\_*.jsonl}, \texttt{run\_state.json}) $\cdot$
Rectángulo con base ondulada = \textbf{documento} (\texttt{Sessions/*.md}) $\cdot$
Rectángulo punteado gris = \textbf{nota explicativa} (sin acción). Flujo: arriba$\rightarrow$abajo, izquierda$\rightarrow$derecha.

\faIcon{tags}~Cada símbolo lleva una \textbf{etiqueta} ({T/P/D/E/A/N} + número, según su tipo); su significado se expande en la tabla \emph{Leyenda de etiquetas} inmediatamente bajo cada diagrama.
\end{tcolorbox}
\end{center}

% ══ Sección 1: Cadena de integridad ═══════════════════════════════════════════
\section{\iconotexto{link}{Cadena de integridad del log append-only}}

El log \texttt{session\_log\_<run>.jsonl} es la fuente de verdad inmutable: cada evento
encadena (\texttt{prev}) el hash SHA-256 de la línea anterior, la secuencia \texttt{seq}
debe ser contigua y la marca \texttt{"append\_only": true} es exigida por
\texttt{integrity}. Cualquier edición, borrado o reordenamiento rompe la cadena.

\begin{center}
\begin{tikzpicture}[flujo, scale=0.72, transform shape]
  \node[almacenamiento] (e1) at (0,0) {A1};
  \node[almacenamiento] (e2) at (4.6,0) {A2};
  \node[almacenamiento] (e3) at (9.2,0) {A3};
  \draw[->] (e1) -- (e2);
  \draw[->] (e2) -- (e3);

  \node[decision] (chk) at (4.6,-2.7) {D1};
  \node[terminador] (ok) at (9.0,-2.7) {T1};
  \node[proceso] (roto) at (0.3,-2.7) {P1};
  \draw[->] (e2.south) -- (chk.north);
  \draw[->] (chk.east) -- node[etiqueta, above]{Sí} (ok.west);
  \draw[->] (chk.west) -- node[etiqueta, above]{No} (roto.east);

  \node[nota] (nota) at (4.6,-5.6) {N1};
  \draw[dashed, grisLinea] (chk.south) -- (nota.north);
\end{tikzpicture}

\vspace{6pt}
\begin{tcolorbox}[cajaContenido, title={\faIcon{table}~Leyenda de etiquetas}]
\small
\begin{tabularx}{\linewidth}{@{}>{\centering\arraybackslash}p{1.4cm}X@{}}
\toprule
\textbf{Etiqueta} & \textbf{Significado} \\ \midrule
A1 & Evento 1: primer evento del log, \texttt{seq=1} y sin \texttt{prev}. \\
A2 & Evento 2: \texttt{seq=2}, \texttt{prev} = hash SHA-256 del evento anterior. \\
A3 & Evento n (último): \texttt{seq=n}, \texttt{prev} = hash del evento n--1 (encadenamiento). \\
D1 & ¿\texttt{integrity} OK? Verifica \texttt{prev}=hash, \texttt{seq} contigua y \texttt{append\_only}. \\
T1 & Log íntegro (tamper-evident): ninguna edición pasó desapercibida. \\
P1 & Manipulación detectable: un desajuste de hash o de secuencia delata edición, borrado o reordenamiento. \\
N1 & Es el primitivo de \emph{hash-chain}, pero sin merkle trees, sin consenso, sin descentralización ni incentivos: tamper-evident sí; blockchain propiamente, no. \\
\bottomrule
\end{tabularx}
\end{tcolorbox}
\end{center}

\faIcon{arrow-right}~\textbf{Contraste con \texttt{run\_state.json}:} el log es
\emph{criptográficamente encadenado} (toda edición es detectable), mientras que la vista
es un JSON plano que se sobrescribe sin rastro. Por eso el log es \textbf{fuente de
verdad} y \texttt{run\_state*.json} una \textbf{vista derivada}.

% ══ Sección 2: Guardia MCP ═══════════════════════════════════════════════════
\section{\iconotexto{shield-alt}{Guardia MCP: ¿el estado leído pertenece a la corrida actual?}}

La guardia usa la comparación de \emph{mtime antes/después} de la corrida. Es la única
señal que vincula temporalmente el archivo al subproceso sin conocer el \texttt{run\_id}
de antemano.

\begin{center}
\begin{tikzpicture}[flujo, scale=0.72, transform shape]
  \node[terminador] (inicio) at (0,0) {T1};
  \node[proceso] (mb) at (0,-1.4) {P1};
  \node[proceso] (run) at (0,-2.9) {P2};
  \node[decision] (existe) at (0,-4.6) {D1};
  \node[terminador] (nada) at (4.8,-4.6) {T2};
  \draw[->] (inicio) -- (mb);
  \draw[->] (mb) -- (run);
  \draw[->] (run) -- (existe);
  \draw[->] (existe.east) -- node[etiqueta, above]{No} (nada.west);

  \node[decision] (nunca) at (0,-6.6) {D2};
  \draw[->] (existe.south) -- node[etiqueta, pos=0.6]{Sí} (nunca.north);

  \node[entradasalida] (leer1) at (0,-8.2) {E1};
  \node[terminador] (fin1) at (0,-9.7) {T3};
  \draw[->] (nunca.south) -- node[etiqueta, pos=0.6]{Sí} (leer1.north);
  \draw[->] (leer1) -- (fin1);

  \node[decision] (mt) at (4.8,-6.6) {D3};
  \draw[->] (nunca.east) -- node[etiqueta, above]{No} (mt.west);

  \node[proceso] (huerfana) at (9.6,-6.6) {P3};
  \draw[->] (mt.east) -- node[etiqueta, above]{No} (huerfana.west);

  \node[entradasalida] (leer2) at (4.8,-8.2) {E2};
  \node[terminador] (fin2) at (4.8,-9.7) {T4};
  \draw[->] (mt.south) -- node[etiqueta, pos=0.6]{Sí} (leer2.north);
  \draw[->] (leer2) -- (fin2);

  \node[nota] (nf) at (9.6,-2.6) {N1};
  \draw[dashed, grisLinea] (run.east) -- (nf.west);
\end{tikzpicture}

\vspace{6pt}
\begin{tcolorbox}[cajaContenido, title={\faIcon{table}~Leyenda de etiquetas}]
\small
\begin{tabularx}{\linewidth}{@{}>{\centering\arraybackslash}p{1.4cm}X@{}}
\toprule
\textbf{Etiqueta} & \textbf{Significado} \\ \midrule
T1 & Corrida MCP lanzada. \\
P1 & Toma \texttt{mtime\_before} = \texttt{mtime(run\_state.json)} antes de la corrida (\texttt{None} si el archivo no existe). \\
P2 & Ejecutar el flujo (\texttt{subprocess}). \\
D1 & ¿Existe \texttt{run\_state.json} al término de la corrida? \\
T2 & \texttt{state\_data = \{\}} (sin vista): el flujo no escribió estado. \\
D2 & ¿\texttt{mtime\_before} = \texttt{None}? (el archivo no existía antes de la corrida). \\
E1 & Leer el estado: fue \emph{creado} durante esta corrida. \\
T3 & Vista pertenece a la corrida. \\
D3 & ¿\texttt{mtime\_ahora} > \texttt{mtime\_before}? (el estado sí fue reescrito en esta corrida). \\
P3 & No usar: vista huérfana (\texttt{mtime} de otra corrida). \\
E2 & Leer el estado: fue \emph{reescrito} durante esta corrida. \\
T4 & Vista pertenece a la corrida. \\
N1 & \texttt{flujo/fin} \textbf{no} sirve de guardia: es un evento de \emph{ciclo de vida} que aparece en los dos casos opuestos (huérfano: el run de hace 1\,h terminó y dejó su estado; legítimo: el flujo recién lanzado terminó bien). Con el criterio ``¿hay flujo/fin?$\rightarrow$descartar'' se descartaría el estado legítimo que el MCP lee \emph{después} de que el flujo termina. \\
\bottomrule
\end{tabularx}
\end{tcolorbox}
\end{center}

% ══ Sección 3: check / clean ═════════════════════════════════════════════════
\section{\iconotexto{stethoscope}{estado\_sesion.py: veredictos y limpieza de la vista}}

Principio conservador: \textbf{``borro solo con certeza''}. \texttt{huerfano}
$\rightarrow$ borra; \texttt{vigente} y \texttt{no\_verificable} $\rightarrow$ conserva;
\texttt{corrupto} $\rightarrow$ conserva y reporta.

\begin{center}
\begin{tikzpicture}[flujo, scale=0.72, transform shape]
  \node[terminador] (inicio) at (0,0) {T1};
  \node[entradasalida] (busca) at (0,-1.4) {E1};
  \node[decision] (existe) at (0,-3.1) {D1};
  \draw[->] (inicio) -- (busca);
  \draw[->] (busca) -- (existe);

  \node[proceso] (nv) at (4.6,-3.1) {P1};
  \node[nota] (nv2) at (4.6,-4.5) {N1};
  \draw[->] (existe.east) -- node[etiqueta, above]{No} (nv.west);
  \draw[dashed, grisLinea] (nv.south) -- (nv2.north);

  \node[proceso] (integ) at (0,-5.2) {P2};
  \draw[->] (existe.south) -- node[etiqueta, pos=0.6]{Sí} (integ.north);

  \node[decision] (integral) at (0,-6.9) {D2};
  \draw[->] (integ) -- (integral);

  \node[proceso] (corrupto) at (4.6,-6.9) {P3};
  \node[nota] (cor2) at (7.8,-6.9) {N2};
  \draw[->] (integral.east) -- node[etiqueta, above]{No} (corrupto.west);
  \draw[dashed, grisLinea] (corrupto.east) -- (cor2.west);

  \node[decision] (ffin) at (0,-8.6) {D3};
  \draw[->] (integral.south) -- node[etiqueta, pos=0.6]{Sí} (ffin.north);

  \node[proceso] (vigente) at (4.6,-8.6) {P4};
  \node[nota] (vig2) at (4.6,-10.0) {N3};
  \draw[->] (ffin.east) -- node[etiqueta, above]{No} (vigente.west);
  \draw[dashed, grisLinea] (vigente.south) -- (vig2.north);

  \node[proceso] (huerfano) at (0,-10.2) {P5};
  \node[nota] (hue2) at (0,-11.6) {N4};
  \node[terminador] (fin) at (0,-13.0) {T2};
  \draw[->] (ffin.south) -- node[etiqueta, pos=0.6]{Sí} (huerfano.north);
  \draw[dashed, grisLinea] (huerfano.south) -- (hue2.north);
  \draw[->] (hue2.south) -- (fin.north);

  \node[nota] (model) at (9.4,-4.4) {N5};
  \draw[dashed, grisLinea] (integ.east) -- (model.west);
\end{tikzpicture}

\vspace{6pt}
\begin{tcolorbox}[cajaContenido, title={\faIcon{table}~Leyenda de etiquetas}]
\small
\begin{tabularx}{\linewidth}{@{}>{\centering\arraybackslash}p{1.4cm}X@{}}
\toprule
\textbf{Etiqueta} & \textbf{Significado} \\ \midrule
T1 & \texttt{check(run\_id)}: inicia la verificación de la vista derivada. \\
E1 & Buscar el log \texttt{session\_log\_<run>.jsonl} del \texttt{run\_id}. \\
D1 & ¿Existe el log para el \texttt{run\_id}? \\
P1 & Veredicto: \texttt{no\_verificable}. \\
N1 & Sin log $\rightarrow$ sin certeza de si la corrida terminó. \texttt{clean} \textbf{NO} actúa (ni lo lista). \\
P2 & \texttt{integrity}: \texttt{prev} = hash $\wedge$ \texttt{seq} contigua $\wedge$ \texttt{append\_only}. \\
D2 & ¿La cadena está íntegra? \\
P3 & Veredicto: \texttt{corrupto}. \\
N2 & Conservar y reportar (incrementa anomalías). \\
D3 & ¿El último evento es \texttt{flujo/fin}? \\
P4 & Veredicto: \texttt{vigente}. \\
N3 & Conservar (en curso, error o vacío). \\
P5 & Veredicto: \texttt{huérfano}. \\
N4 & Único caso que \texttt{clean} borra (\texttt{unlink}). \\
T2 & Vista eliminada. \\
N5 & \texttt{modelo\_obsoleto}: solo \emph{advertencia} informativa (\texttt{deepseek-v4-pro}); \textbf{no} altera veredicto ni incrementa anomalías. \\
\bottomrule
\end{tabularx}
\end{tcolorbox}
\end{center}

El modelo obsoleto es una señal de atención (``reanudar con este modelo fallaría''), no
un criterio de limpieza: en el caso real detectado se borró por \texttt{huerfano}, no por
el modelo.

% ══ Sección 4: bitacoras.py ══════════════════════════════════════════════════
\section{\iconotexto{book}{Eslabón semántico: bitacoras.py (\texttt{nueva} y \texttt{check})}}

El subcomando \texttt{nueva} crea la bitácora canónica (nunca sobrescribe); \texttt{check}
clasifica cada bitácora por su fecha relativa a \texttt{FECHA\_PLANTILLA} (2026-09-11).

\begin{center}
\begin{tikzpicture}[flujo, scale=0.72, transform shape]
  % Columna A: nueva
  \node[terminador] (nt) at (0,0) {T1};
  \node[proceso] (n1) at (0,-1.5) {P1};
  \node[decision] (n2) at (0,-3.3) {D1};
  \draw[->] (nt) -- (n1);
  \draw[->] (n1) -- (n2);

  \node[terminador] (n3) at (4.2,-3.3) {T2};
  \draw[->] (n2.east) -- node[etiqueta, above]{Sí} (n3.west);

  \node[proceso] (n4) at (0,-5.6) {P2};
  \draw[->] (n2.south) -- node[etiqueta, pos=0.6]{No} (n4.north);
  \node[terminador] (n5) at (0,-7.4) {T3};
  \draw[->] (n4) -- (n5);

  % Columna B: check
  \node[terminador] (ct) at (7.4,0) {T4};
  \node[entradasalida] (c1) at (7.4,-1.5) {E1};
  \node[decision] (c2) at (7.4,-3.3) {D2};
  \draw[->] (ct) -- (c1);
  \draw[->] (c1) -- (c2);

  \node[proceso] (c3) at (11.2,-3.3) {P3};
  \draw[->] (c2.east) -- node[etiqueta, above]{Sí} (c3.west);

  \node[decision] (c4) at (7.4,-5.0) {D3};
  \draw[->] (c2.south) -- node[etiqueta, pos=0.6]{No} (c4.north);

  \node[proceso] (c5) at (11.2,-5.3) {P4};
  \draw[->] (c4.east) -- node[etiqueta, above]{No} (c5.west);

  \node[decision] (c6) at (7.4,-6.8) {D4};
  \draw[->] (c4.south) -- node[etiqueta, pos=0.6]{Sí} (c6.north);

  \node[proceso] (c7) at (11.2,-6.8) {P5};
  \draw[->] (c6.east) -- node[etiqueta, above]{No} (c7.west);

  \node[proceso] (c8) at (7.4,-8.6) {P6};
  \draw[->] (c6.south) -- node[etiqueta, pos=0.6]{Sí} (c8.north);

  \node[nota] (nleg) at (7.4,-10.6) {N1};
  \draw[dashed, grisLinea] (c8.south) -- (nleg.north);
\end{tikzpicture}

\vspace{6pt}
\begin{tcolorbox}[cajaContenido, title={\faIcon{table}~Leyenda de etiquetas}]
\small
\begin{tabularx}{\linewidth}{@{}>{\centering\arraybackslash}p{1.4cm}X@{}}
\toprule
\textbf{Etiqueta} & \textbf{Significado} \\ \midrule
\multicolumn{2}{@{}l}{\textit{Columna A --- ``nueva''}} \\ \midrule
T1 & \texttt{nueva --tema}: crea la bitácora canónica de la sesión. \\
P1 & Slug ASCII del tema $\rightarrow$ archivo \texttt{Sessions/YYYY-MM-DD\_<slug>.md}. \\
D1 & ¿Ya existe el archivo? \\
T2 & \texttt{ya\_existia} (nunca sobrescribe). \\
P2 & Escribir las 5 secciones: Tema $\cdot$ Contexto $\cdot$ Decisiones (usuario) $\cdot$ Actividades $\cdot$ Pendientes. \\
T3 & Lista para \texttt{check}. \\ \midrule
\multicolumn{2}{@{}l}{\textit{Columna B --- ``check''}} \\ \midrule
T4 & \texttt{check}: clasifica cada bitácora según su fecha relativa a \texttt{FECHA\_PLANTILLA}. \\
E1 & Listar \texttt{Sessions/*.md}. \\
D2 & ¿Fecha del nombre < \texttt{FECHA\_PLANTILLA} (2026-09-11)? \\
P3 & Legado: \texttt{ok} forzado \texttt{True} (informativo, no se reescribe). \\
D3 & ¿Secciones completas (hoy/reciente)? \\
P4 & Anomalía accionable: afecta veredicto y \texttt{exit code}. \\
D4 & ¿\texttt{hay\_bitacora} para hoy? \\
P5 & Veredicto: ``atencion'' (día sin bitácora). \\
P6 & Veredicto: \texttt{ok}. \\
N1 & Legado (antes de 2026-09-11): escritas con otra convención (ej. las 33 antiguas sin \texttt{Decisiones (usuario)}). Reescribirlas sería falsificar historia. \\
\bottomrule
\end{tabularx}
\end{tcolorbox}
\end{center}

En la columna derecha, ``hoy'' o ``reciente'' sin secciones completas es una anomalía
real a corregir \emph{antes} de cerrar; la salida cruza con \texttt{sesion\_hoy},
booleano que \texttt{check} obtiene filtrando por el prefijo del nombre
(\texttt{startswith("YYYY-MM-DD")}).

% ══ Sección 5: cierre de sesión ══════════════════════════════════════════════
\section{\iconotexto{door-closed}{Cierre de sesión: remediación obligatoria (edge case de la directiva)}}

Si al cierre \texttt{check} reporta \emph{falta de bitácora del día} o \emph{secciones
vacías} (hoy/reciente), la sesión no debe cerrar así: la recuperación es correctiva y el
defecto no es cosmético, porque el contexto semántico es irrecuperable una vez cerrada la
sesión (no se regenera desde ningún log).

\begin{center}
\begin{tikzpicture}[flujo, scale=0.72, transform shape]
  \node[terminador] (inicio) at (0,0) {T1};
  \node[proceso] (p1) at (0,-1.5) {P1};
  \node[decision] (d1) at (0,-3.3) {D1};
  \draw[->] (inicio) -- (p1);
  \draw[->] (p1) -- (d1);

  \node[terminador] (fin) at (0,-5.2) {T2};
  \draw[->] (d1.south) -- node[etiqueta, pos=0.7]{No} (fin.north);

  \node[proceso] (r1) at (5.6,-3.3) {P2};
  \node[proceso] (r2) at (5.6,-4.6) {P3};
  \node[proceso] (r3) at (5.6,-5.9) {P4};
  \draw[->] (d1.east) -- node[etiqueta, above]{Sí} (r1.west);
  \draw[->] (r1) -- (r2);
  \draw[->] (r2) -- (r3);
  \draw[->] (r3.west) -- ++(-1.8,0) |- ([xshift=8pt]d1.east) node[etiqueta, pos=0.35, right=6pt]{revalidar};

  \node[nota] (nleg) at (5.6,-8.6) {N1};
\end{tikzpicture}

\vspace{6pt}
\begin{tcolorbox}[cajaContenido, title={\faIcon{table}~Leyenda de etiquetas}]
\small
\begin{tabularx}{\linewidth}{@{}>{\centering\arraybackslash}p{1.4cm}X@{}}
\toprule
\textbf{Etiqueta} & \textbf{Significado} \\ \midrule
T1 & Cierre de sesión. \\
P1 & Ejecutar \texttt{bitacoras.py check}: verifica la bitácora del día y el llenado de secciones. \\
D1 & ¿Falta la bitácora del día o hay secciones vacías (hoy/reciente)? \\
T2 & Cerrar sesión (sin defectos de bitácora). \\
P2 & \texttt{nueva --tema} (si falta la bitácora del día). \\
P3 & Rellenar las secciones semánticas. \\
P4 & Revalidar con \texttt{check}. \\
N1 & \textbf{Exento:} las bitácoras legado (antes de 2026-09-11) son informativas y no se reescriben. \\
\bottomrule
\end{tabularx}
\end{tcolorbox}
\end{center}

\begin{cajaConcepto}
\faIcon{lightbulb}~\textbf{Información que reside \emph{exclusivamente} en \texttt{Sessions/}:}
el \textbf{contexto semántico}---el POR QUÉ de las decisiones (sobre todo las del usuario),
lo que se descartó y por qué, intenciones y matices, reglas y aprendizajes derivados, y
pendientes acordados con su contexto. Bajo nivel = datos reproducibles (log + run\_state);
alto nivel = significado (\texttt{Sessions/}). Sin bitácora, esa información se pierde
\emph{definitivamente}.
\end{cajaConcepto}

\end{document}